\scanpage{53}
\mspage{5}
% archive foliation: 52
\[
 a)\longrightarrow b)\qquad \Gamma=\underline{\operatorname{Norm}}_G\underline V
\]
$b)\longrightarrow a)$\quad $\underline V$ est l'unique sous-schéma en groupes \inserted{distingué} fini étale de $\Gamma$, \deleted{\unsure{q}} de degré $4$, \deleted{distingué}.
\margin{\uline{NB}\quad $\Gamma/V\simeq\underline{\operatorname{Aut}}_{S\mathrm{gr}}(V)\simeq\underline{\operatorname{Aut}}_S(V^*)$}
\[
 c)\longrightarrow b)\qquad
 \Gamma=\underline{\operatorname{Aut}}_S(X,\Delta)\xleftarrow{\sim}\underline{\operatorname{Aut}}^+_{S\text{-}\mathrm{oct}}(\Delta)
\]
$\unsure{a})\longrightarrow c)$\quad \deleted{$\underline V$ est} Toute section $\inserted{\sigma\in}V^*=V\setminus(\text{section unité})$ définit une action \inserted{fidèle} de \deleted{\unsure{ce}} $(\mathbb Z/2\mathbb Z)_S$ sur $X$, donc un ``axe'' $\Delta_\sigma$ de $X$, et $\Delta$ est \deleted{la réuni\unsure{on}} (loc. ét.) le \deleted{\unsure{sché}} diviseur somme de ses axes\ldots

\uline{Corollaire 1} Pour un diviseur étale de degré $6$ $\Delta$ dans $X$, il existe au plus une structure icosaédrale régulière ; \deleted{dont le groupe des autom.} \inserted{S'il en existe une, le} \inserted{schéma en} groupe de ses automorphismes \inserted{directs} est isom.
\[
\underline{\operatorname{Aut}}_S^{\unsure{+}}(X,\Delta).
\]

\uline{Corollaire 2}\quad La catégorie des couples $(X,\Delta)$ \inserted{munis d'un iso $\underline{\operatorname{Or}}\Delta/S\simeq\mu_{4S}^*$ (on dit qu'ils sont $\mu_4^*$-orientés)} sur $S$, (i.e. $(X,\underline V)$, \deleted{\unsure{orientés}} $(X,\underline\Gamma)$, \inserted{$\simeq$ \deleted{\unsure{$V$}} octaèdres combinatoires $\Delta\ldots$}) est équivalente à celle des \deleted{\unsure{sur}} $S$, ou encore : celle des revêtements finis étales \deleted{$\unsure{\Delta}$} \inserted{$\widetilde F$} de degré $4$ de $S$, en associant à \deleted{$\unsure{H}$} $(X,\Delta)$ le revêtt des quatre ``faces doubles antipodiques'' de la structure octaédrale $\Delta$, ou ce qui revient au $=$, celui des sous-$3$-groupes de Sylow de \deleted{$\unsure{\mathfrak S_4}$}
\[
\underline{\operatorname{Aut}}^+_{S\mathrm{oct}}(\Delta)\simeq\underline{\operatorname{Aut}}_S(X,\Delta).
\]
La structure-type${}^*$ correspond \deleted{\unsure{au choix de}} $X\subset\mathbb P^2_{\mathbb Z[\frac12]}$, $X=V(x^2+y^2+z^2=0)$, avec l'op. de \deleted{\unsure{$G$}} $\{\pm\}^3/\mathrm{diag}\simeq V_{\{0,1,\infty\}}(\mathbb F_2)$ donnée par les chgts de signes en $x,y,z$${}^*$.

\margin{${}^*$ La \unsure{biface|face} distinguée est formée des deux faces $\{(0,1,i),(i,0,1),(1,i,0)\}$ et $\{(0,1,-i),(-i,0,1),(1,-i,0)\}$.}
\margin{\uline{NB} Ici $X$ est totalement anisotrope sur $\mathbb R$ i.e. n'a pas de pts réels sur $\mathbb R$ -- le revêtement galoisien associé étant de groupe $\unsure{V}_{\{0,1,\infty\}}(\mathbb F_2)$ de $\mathbb P^1_{\unsure{\mathbb C}}$, ramifié \unsure{seulement} en $0,1,\infty$\ldots}
